\dresumoheader{mieiro}{Manoel Jorge de Lima}{Attention and Engagement in Remote Classes: The Challenge of the Expository Method}{2025}{Undergraduate Thesis}{Federal Center for Technological Education Celso Suckow da Fonseca}{Rio de Janeiro}{Rio de Janeiro}{2025}

This study aims to analyze the degree of distraction and low engagement of classes subjected to expository lessons in a remote learning environment. To this end, a web browser extension called ESPEON was developed, which actively monitors students’ behavior during classes, recording navigation data and the use of input peripherals (camera and microphone). The research was conducted with classes of five students, segmented by field of knowledge, exposed to expository lessons delivered by instructors from the Centro Federal de Educação Tecnológica Celso Suckow da Fonseca (CEFET/RJ), via the Microsoft Teams platform. At the end of each class, the tool generated reports with indicators and metrics indicative of students’ attention and interactivity. The quantitative data obtained were analyzed using descriptive statistics to evaluate the performance of the expository methodology in the virtual classroom.